\section{异构性才是这里的 Scale}

互联网基础设施常用请求数、节点数或数据量描述 scale。Palantir 面对的困难还有另一维：同一软件必须跨越完全不同的网络、硬件、监管、更新窗口和故障模式。一个大型云集群可能有统一控制面，数百或上千个客户环境却各有边界。系统的扩展单位因此不是单台服务器，也不只是一个 Kubernetes cluster，而是约束组合的数量。

\subsection{从 On-Prem Monolith 到自动交付}

Palantir 创业早期，AWS S3 尚未成为可依赖的公共基础设施，客户系统大量运行在 on-premises 环境。\term{On-prem} 指软件部署在客户自有机房或设备，而非供应商控制的公共云。很多政府网络还是 \term{air-gapped}，即与普通互联网物理或逻辑隔离，更新包和诊断信息需要受控转运。云时代的远程 root access、即时扩容与统一镜像在这里都不能假设存在。

\begin{figure}[H]
\centering
\includegraphics[width=0.94\textwidth]{images/01-infrastructure-epochs.png}
\caption{访谈中的基础设施演进：手工 on-prem、scale-out、microservices，再到 Apollo 与 Rubix。}
\figsource{依据字幕 05:00--13:30 重绘，`images/01-infrastructure-epochs.png`。}
\end{figure}

\begin{knowledgebox}{读图：每次抽象都会暴露新的控制问题}
Monolith 组件少，却难以独立发布；microservices 让团队并行开发，却增加版本、schema、依赖和部署顺序。Kubernetes 统一容器调度，但不能自动理解每个软件的升级前置条件，也不能替组织定义安全策略。Apollo/Rubix 并非“再加两层”，而是把前一阶段留给人工记忆的控制逻辑编码成平台能力。
\end{knowledgebox}

\begin{importantbox}{术语消化：Microservice 与 SRE}
Microservice 是可独立开发、发布和扩缩容的服务单元，服务之间通过 API 或消息协作。SRE 是 Site Reliability Engineer，站点可靠性工程师，负责可用性、容量、事故响应和自动化。早期把 SRE 直接派到每个客户现场可以解决问题，却把组织人数变成扩展上限。
\end{importantbox}

\subsection{四类环境差异}

讲者用粗略数量级描述课堂时点的内部规模：约千个生产环境、其中约百个断网环境、数千个服务和大量每周升级。这里应把数字看作口述的规模感，而不是财务披露。更稳定的结论是环境跨度：有的运行在大型多租户云，有的在客户单租户集群，有的靠受控离线包更新，有的位于战术边缘或资源受限设备。代码必须面对不同连接性、容量和认证规则。

\begin{figure}[H]
\centering
\includegraphics[width=0.94\textwidth]{images/02-heterogeneous-fleet.png}
\caption{异构 fleet 的四个代表性端点：断网、边缘、单租户与多租户云。}
\figsource{依据字幕 05:00--09:00 重绘，`images/02-heterogeneous-fleet.png`。}
\end{figure}

\begin{knowledgebox}{读图：环境描述必须是多维向量}
只写“运行在云上”不足以决定部署策略。至少要记录连接模式、硬件和架构、可用容量、合规级别、数据驻留、维护窗口、允许的操作人员和故障恢复方式。环境 metadata 越完整，控制面越能自动计划；信息缺失时，自动化只会把错误更快传播到 fleet。
\end{knowledgebox}

\begin{table}[H]
\centering
\begin{tabular}{L{0.19\textwidth}L{0.31\textwidth}L{0.4\textwidth}}
\toprule
维度 & 典型差异 & 对交付系统的影响 \\
\midrule
连接性 & 在线、间歇、单向、完全断网 & 计划同步、artifact 传输与遥测延迟 \\
算力 & 大型云、固定机房、边缘设备 & 资源 request、扩缩容与功能裁剪 \\
安全 & 商业基线、监管、分类网络 & 身份、审批、审计和 operator access \\
拓扑 & 单租户、多租户、跨区域 & 隔离、容量池和故障域 \\
变更窗口 & 连续发布、定期窗口、人工接入 & release channel、bundle 与回滚策略 \\
\bottomrule
\end{tabular}
\caption{异构环境不是部署标签，而是一组会改变控制算法的约束。}
\end{table}

\begin{warningbox}{规模数字需要绑定时间和口径}
Palantir 2022 Apollo Demo Day 官方材料写的是 250+ distinct customer environments，课堂口述使用了更高的粗略 production-environment 数量级。二者可能来自不同年份、产品范围或计数口径。讲义保留两处来源，不用一个数字证明另一个数字错误，也不把口述近似写成当前官方统计。
\end{warningbox}

\subsection{本章小结}

Palantir 的 scale 主要表现为 fleet 异构性与变更频率。Air gap、边缘资源、单租户政策和多租户云共同要求软件把环境约束建模，而不是假设所有机器属于同一控制域。Apollo 的出现，就是为了让这种环境模型参与每次升级决策。

